\chapter{Теоретична основа}
\label{ch:osnova}

Тази глава въвежда понятията, необходими за разбирането на останалата част от
настоящата дипломна работа: ролята на динамиката във възприятието на ударните инструменти,
представянето на такива партии в MIDI, признаците, които могат да се извлекат
от структурата и времето на партията, както и набора от данни E-GMD, върху
който се провеждат всички експерименти.

\section{Динамика и възприятие на ударните инструменти}
\label{sec:dinamika}

Под \emph{динамика} на една нота разбираме силата на удара, кодирана в MIDI като
\eng{velocity} — цяло число от $0$ до $127$. За разлика от много мелодични
инструменти, при барабаните \eng{velocity} влияе не само на силата на звука, а и
на \emph{тембъра}: тихо ударен барабан звучи качествено различно от силно ударен
(при акустичните барабани заради нелинейността на удара, а при електронните и
заради превключване на семпли в зависимост от силата, вж.\ раздел~\ref{sec:egmd}).

Три явления правят динамиката особено важна за ударните инструменти. Първо, \emph{акцентите}
— по-силните удари, които маркират метричната структура и създават усещане за
ритъм. Второ, \emph{призрачните ноти} (\eng{ghost notes}) — много тихи удари
(най-често по соло барабана), които създават пълнеж, без да се възприемат като
самостоятелни акценти. Трето, постепенните засилвания и затихвания (т.нар. \emph{крешендо}
и \emph{декрешендо}), както и стилистичните конвенции (напр.\ акценти на слабо време във 
фънка и др. жанрове). Разпределението на \eng{velocity} за много видове ударни инструменти
следователно е \emph{многомодално} — например при соло барабана ясно се отделят два режима
(призрачни ноти и нормални/акцентирани удари), както се вижда по-нататък на
фиг.~\ref{fig:velocity_dists}. Този факт има пряко следствие за моделирането:
точкова оценка, която връща една стойност на нота, неизбежно ще „осреднява'' и
ще заличава именно тази вариация, която носи човешкото усещане.

\section{MIDI и представяне на барабанни партии}
\label{sec:midi}

MIDI (Musical Instrument Digital Interface) е протокол, който представя музиката 
като поток от събития, а не като звуков сигнал. За целите на тази работа 
всяка барабанна нота се описва като кортеж, състоящ се от \emph{начално време} (\eng{onset}), 
\emph{продължителност} (\eng{duration}), \emph{височина} (\eng{pitch}, цяло число $ \in [0, 127] $),
\emph{канал} ($ 1- 16 $) и \emph{сила} (\eng{velocity}, цяло число $ \in [0, 127] $).
Барабаните по стандарт заемат канал~10 (индекс~9), където, по спецификацията
на \eng{General MIDI}~\cite{gm1991spec}, всяка различна стойност за височина
съответства на определен ударен инструмент — напр.\ $36$ е бас каса, $38$ —
акустичен соло барабан, $42$ — затворен фус, и т.н.

\eng{General MIDI} (за краткост — GM) обаче дефинира недостатъчно на брой перкусионни слотове
за целите на електронни барабани и падове, имитиращи звуковото разнообразие и артикулации на истинските
акустични комплекти. Затова някои производители (напр.\ Roland) използват разширени и понякога
несъвместими с GM разпределения за модулите си. Това налага създаването в насотящата работа на
нормализация, която категоризира дефинираните от GM височини в \emph{канонични гласове} (\eng{voices}) — напр.\
всички варианти на удари по фус се групират в семейство фус гласове.
Дефинирането на това групиране не бива да се прави произволно; в
раздел~\ref{sec:egmd} и в глава~\ref{ch:metod} то се извежда по начин, ръководен от данните.

\section{Признаци от структура и време}
\label{sec:priznaci}

Централно ограничение на задачата е, че при извод входът е „плоско'' MIDI, чиито
\eng{velocity} стойности са нарочно премахнати и тепърва трябва да се изградят.
Целият контекст трябва да се изгради само от музикална структура и време:
\emph{кой} инструмент свири и \emph{кога}.

Музиката се харектеризира с периодичност, която се проявява на две различни нива:
\emph{такт} (\eng{bar} или \eng{measure}) и \emph{време} (\eng{beat}).
Един такт е обикновено се възприема като най-малката част, която се повтаря ритмично, а често и
хармонично или мелодично. Един такт се дели на няколко времена, в зависимост от \emph{размера}
(\eng{time signature}) на произведението. Най-често срещаният размер в музиката е $4/4$, който ни казва, че тактовете
се делят на четири равни времена, всяко от които бележим чрез четвъртина нота в писмен вид.
Времето (\eng{beat}) е и най-малката единица, чрез която измерваме темпо в музиката. Величината,
с която отмерваме точно темпо, наричаме \emph{bpm} (\eng{beats per minute}). Тя ни казва колко времена се съдържат в една минута,
което е мярка за скоростта на изпълнение.

Ключово понятие в настоящата работа е \emph{метричната фаза} — позицията на нотата спрямо
времето и такта, изразена като реално число в $[0,1)$. Фазата е \emph{циклична}
величина: стойности $0{,}99$ и $0{,}01$ са музикално съседни, но числено далечни.
Това различие би се отразило зле при създаване на модели, особено за невронни мрежи
по време на градиентното спускане при обучението им.
Затова позицията спрямо такт или време се кодира върху единичната окръжност
чрез двойката $[\sin(2\pi\varphi),\cos(2\pi\varphi)]$, което премахва изкуствения скок на
границата на времето.

За разлика от фазата, \emph{времевите разстояния} между
удари (напр.\ време до предишен или следващ удар) измерват абсолютни интервали, а
не позиции спрямо периодични явления, и се кодират като логаритмично мащабирани скалари в
единици „време''. Освен това, като входни данни от структурата се извличат локална \emph{плътност}
на ударите, \emph{едновременност} (кои гласове звучат в един и същ момент) и
\emph{стилови} признаци (жанр, темпо). Обосновката и пълният списък на признаците
са дадени в глава~\ref{ch:metod}.

\section{Наборът от данни E-GMD}
\label{sec:egmd}

\eng{Expanded Groove MIDI Dataset} (E-GMD)~\cite{callender2020egmd} е разширение на
\eng{Groove MIDI Dataset} и, доколкото ни е известно, първият голям набор с
\emph{човешки изпълнения}, които съдържат в себе си анотации и на \eng{ MIDI velocity} за барабани.
Съдържа изпълнения от няколко барабаниста, записани на електронен комплект Roland,
с предварително направено разделяне на подмножества за обучение, валидация и тест
(съответно $35\,217$ / $5\,031$ / $5\,289$ MIDI файла). Всеки файл съдържа едно
изпълнение; метаданните включват стил (жанр), темпо (\eng{bpm}), тактов
размер и тип на изпълнението (цяла фраза или само брейк).
Изследователски анализ върху подмножеството за обучение (11{,}07 млн.\ ноти,
вж.\ глава~\ref{ch:metod}) установи няколко
факта, съществени за метода:

\begin{itemize}
  \item \textbf{Тактовият размер е практически константа:} $98{,}78\%$ от файловете
        са в $4/4$, а неравноделните размери ($5/4$, $5/8$) са под
        $0{,}3\%$. Това позволява тактовият размер да се пренебрегне като признак
        в последователния модел.
  \item \textbf{Появяват се 25 GM височини}, две от които са \emph{извън}
        \eng{General MIDI} — артикулации по „ръба'' (\eng{edge}) на фуса на
        Roland комплекта, използван за записване на набора от данни.
        Разпределението на употребата и на \eng{velocity} по
        височини (фиг.~\ref{fig:pitch_counts} и~\ref{fig:velocity_dists}) е
        основата за каноничното групиране на гласовете.
  \item \textbf{Соло барабанът има висока дисперсия на силата на удар} ($\sigma\approx36$, обхват от
        призрачни ноти до акценти) — пряка мотивация за вероятностните модели в
        глава~\ref{ch:metod}.
\end{itemize}

\begin{figure}[H]
  \centering
  \includegraphics[width=0.9\textwidth]{fig_velocity_dists.png}
  \caption{Разпределения на \eng{velocity} по канонични гласове (в обучаващото
  подмножество на E-GMD). Ясно се вижда многомодалността на соло барабана
  (призрачни ноти и акценти) и почти винаги силните удари по ,,електронния" соло барабан
  (обозначен като \eng{Electric Snare}) и тези по ръба на фуса (\eng{Hi-Hat Edge}).}
  \label{fig:velocity_dists}
\end{figure}

\begin{figure}[H]
  \centering
  \includegraphics[width=0.85\textwidth]{fig_pitch_counts.png}
  \caption{Хистограма на GM височините в обучаващото подмножество.
  Разпределението е с ,,дълга опашка": няколко гласа (бас каса, соло барабан, фус,
  райд) доминират, а много перкусионни слотове са редки или изобщо липсват.}
  \label{fig:pitch_counts}
\end{figure}

\paragraph{Артефакт от записване през множество комплекти.}
В E-GMD, всяко изпълнение е повторно записано на 43 различни бленди (\eng{kits}) на модула
Roland. Оказва се, че част от блендите \emph{пренасочват} падовете към различни
MIDI височини (напр.\ первазен удар по соло барабана понякога присъства така, но понякога е
отчетен просто като соло барабан), докато времето и силата остават почти непроменени.
Тъй като силата (\eng{velocity}) — целевата ни величина — е практически инвариантна спрямо
блендата, това не засяга самата задача за предсказване на динамиката, но внася шум в признака
„глас'' (MIDI височина). Явлението и неговите последствия се разглеждат по-подробно във връзка с
анализа на грешките (раздел~\ref{sec:interpretiruemost}) и сред ограниченията на
работата (раздел~\ref{sec:budeshti}).
